\setupcolors[state=start]
\setupcolor[hex]
\definecolor[darkblue][h=003366]
\definecolor[darkgray][h=e6e6e6]
\definecolor[redline][h=d9534f]
\definecolor[blueline][h=5bc0de]
\definecolor[greenline][h=5cb85c]

\setuphead[section][style=\bfa, color=darkblue, before=\blank, after=\blank]
\setuphead[subsection][style=\bf, color=black, before=\blank, after=\blank]
\setuppapersize[A4][A4]
\setupbodyfont[11pt]
\setuplayout[width=middle, height=middle, topspace=2cm, bottomspace=2cm]

% Agregados a la cabecera de la guía de estilo del repositorio de implementación
% (docs/guides/context_style_guide.md): idioma, sub-subsecciones con la misma regla que
% las subsecciones, índice de secciones y subsecciones, y rutas del código en
% monoespaciado con corte de línea en / y _.
\mainlanguage[es]
\setuphead[subsubsection][style=\bf, color=black, before=\blank, after=\blank]
\setupcombinedlist[content][list={section,subsection}]
\define[1]\ruta{{\tt\hyphenatedurl{#1}}}

\starttext

\title{\bf Validación de los modelos OPP v5:\crlf equivalencia entre el monolítico y el Split}
\blank[big]

{\it Notas de trabajo para la tesina. Describen cómo se verificó que el modelo monolítico
(\ruta{opp_v5_1_mono}) y el modelo en dos etapas (\ruta{opp_v5_2_split}) son equivalentes: qué
significa la equivalencia, con qué métodos se comprobó, con qué criterio se aceptó cada
comprobación, qué resultados dio cada una y qué queda sin cubrir. Complementan las notas de
desarrollo (\ruta{Notas/opp_v5_development.tex}), que tienen la formulación completa de los dos
modelos; acá se repite solo lo necesario. Todos los números salen de los archivos de resultados
del repositorio de implementación (\ruta{linear-picking-routing}) y se volvieron a calcular desde
esos archivos al escribir estas notas.}

\blank[big]

\completecontent

\page

%%%%%%%%%%%%%%%%%%%%%%%%%%%%%%%%%%%%%%%%%%%%%%%%%%%%%%%%%%%%%%%%%%%%
\startsection[title={Qué se quiere validar}, reference=sec:objetivo]

El Split se usa en la tesina como la referencia exacta del problema: contra su óptimo se miden
las heurísticas, y en los estudios de escala es la referencia contra la que se mide al
monolítico. Eso solo es legítimo si el Split resuelve {\em el mismo problema} que el
monolítico, que es la formulación directa del problema. La validación responde esa pregunta.

\startsubsection[title={Definición de equivalencia}]

Se dice que los dos modelos son equivalentes si, para cualquier instancia:

\startitemize[n]
\item {\bf Mismo estado:} uno es factible si y solo si el otro lo es.
\item {\bf Mismo valor óptimo:} cuando los dos prueban optimalidad, el objetivo es el mismo,
salvo la tolerancia numérica del solver.
\stopitemize

{\bf No se exige la misma solución.} Los dos modelos pueden devolver soluciones óptimas
distintas (otras rutas, otra asignación de lotes a pickers) con el mismo objetivo. Esa diferencia
no es un error: se clasifica como {\em óptimo alternativo} y se explica (sección
\in[sec:alternativos]). Separar el objetivo de la solución es una decisión metodológica central:
un test que exigiera una solución puntual fallaría en un modelo correcto.

{\bf Tolerancia numérica.} Se considera que dos objetivos $a$ y $b$ son iguales si
\startformula
|a - b| \le 10^{-6} \cdot \max(1, |a|, |b|),
\stopformula
que es del orden de las tolerancias por defecto de SCIP. Las diferencias que superan ese umbral
se reportan aparte, con su explicación (sección \in[sec:pares]).
\stopsubsection

\startsubsection[title={Por qué no alcanza con mirar que los tests pasen}]

Que los dos modelos pasen la misma suite de tests muestra que cada uno respeta las reglas que
los tests verifican, pero no que den el mismo óptimo: muchos tests verifican propiedades de la
solución (cobertura de ubicaciones, sin solapamiento, asignación de cada lote) y no su valor. Por
eso, además de los tests, se comparan los dos modelos entre sí, instancia por instancia, sobre
las soluciones guardadas.
\stopsubsection

\stopsection

%%%%%%%%%%%%%%%%%%%%%%%%%%%%%%%%%%%%%%%%%%%%%%%%%%%%%%%%%%%%%%%%%%%%
\startsection[title={Metodología}, reference=sec:metodologia]

La validación combina cinco líneas de evidencia. Cada una responde una pregunta distinta y cubre
lo que las anteriores no alcanzan:

\blank[medium]
\bTABLE[split=yes, align=middle, frame=on, width=broad]
\setupTABLE[row][1][background=color, backgroundcolor=darkgray, style=bold]
\setupTABLE[column][2][align=flushleft]
\setupTABLE[column][3][align=flushleft]
\setupTABLE[column][4][align=flushleft]
\bTR \bTH Línea \eTH \bTH Pregunta \eTH \bTH Instancias \eTH \bTH Criterio de aceptación \eTH \eTR
\bTR \bTD 1. Argumento formal \eTD \bTD ¿Deben coincidir los dos modelos? ¿Bajo qué condiciones? \eTD \bTD Todas \eTD \bTD Demostración en los dos sentidos \eTD \eTR
\bTR \bTD 2. Suite de tests compartida \eTD \bTD ¿Cada modelo cumple las reglas y da los valores calculados a mano? \eTD \bTD 64 tests funcionales (63 en el monolítico) y 19 de stress \eTD \bTD Aserciones de cada test \eTD \eTR
\bTR \bTD 3. Comparación cruzada de soluciones \eTD \bTD ¿Mismo estado y mismo objetivo en cada instancia? ¿Misma solución? \eTD \bTD Todas las instancias funcionales y de stress \eTD \bTD Estado igual y objetivo igual con tolerancia relativa $10^{-6}$ \eTD \eTR
\bTR \bTD 4. Verificación cruzada de factibilidad \eTD \bTD Donde el monolítico no resuelve, ¿la solución del Split es factible para él? \eTD \bTD Los 9 stress que resuelve el Split \eTD \bTD Veredicto \quote{compliant} del verificador \eTD \eTR
\bTR \bTD 5. Pares de la exploración \eTD \bTD ¿Se mantiene en instancias más grandes y realistas? \eTD \bTD 71 pares de los estudios 08 a 11 \eTD \bTD Óptimo probado por los dos; tolerancia relativa $10^{-6}$ \eTD \eTR
\eTABLE
\blank[medium]

\startsubsection[title={Principios}]
\startitemize
\item {\bf Misma instancia por construcción:} los dos modelos reciben la instancia de la misma
llamada al mismo constructor de escenarios, y los dos escriben sus datos con el mismo módulo
(\ruta{build_dat_files}). Ninguna instancia se reconstruye a mano para un modelo.
\item {\bf Comparar sin volver a resolver:} cada corrida guarda su solución completa; la
comparación lee esos archivos. Es barata, se puede repetir y no depende del reloj ni de la
máquina.
\item {\bf Objetivo y solución por separado:} la equivalencia se juzga por estado y objetivo; la
igualdad de soluciones se mide aparte y cada diferencia se clasifica.
\item {\bf Un modelo verifica al otro sin reescribirlo:} para comprobar una solución contra el
monolítico, se fijan sus valores sobre las variables del monolítico compilado tal como está; no
se copia ni se reescribe ninguna restricción.
\item {\bf Tolerancias explícitas:} cada comparación declara su tolerancia, y las diferencias que
la superan se informan aunque sean mínimas.
\item {\bf Toda discrepancia se explica o queda registrada como pendiente:} ninguna diferencia se
descarta como ruido sin una explicación o una hipótesis declarada.
\stopitemize
\stopsubsection

\startsubsection[title={Orden lógico}]

El argumento formal (línea 1) dice qué esperar y bajo qué condiciones. La suite de tests (línea 2)
valida cada modelo por separado contra valores conocidos. La comparación cruzada (línea 3)
contrasta los dos modelos instancia por instancia, pero solo donde los dos terminan, que en la
práctica son las instancias funcionales. En stress el monolítico no termina ninguna instancia,
así que la línea 4 cambia la pregunta: si no se puede comparar el óptimo, al menos se comprueba
que lo que devuelve el Split sea factible y tenga el mismo costo bajo las restricciones del
monolítico. La línea 5 extiende la comparación de objetivos a instancias del perfil realista,
más grandes que las funcionales y todavía al alcance del monolítico.
\stopsubsection

\stopsection

%%%%%%%%%%%%%%%%%%%%%%%%%%%%%%%%%%%%%%%%%%%%%%%%%%%%%%%%%%%%%%%%%%%%
\startsection[title={Línea 1: el argumento formal}, reference=sec:formal]

En el repositorio de implementación el argumento estaba escrito en una sola frase: la ruta más
corta domina a cualquier otra, porque el objetivo no tiene un término de ruta y una ruta más
corta solo relaja las restricciones de duración y de cadena de frío. Acá se escribe completo, en
los dos sentidos, a partir de las dos formulaciones, y con las condiciones que necesita.

\startsubsection[title={Las dos formulaciones, en resumen}]

Se usa la notación de las notas de desarrollo. El monolítico tiene las restricciones (1) a (22).
Para no mezclar numeraciones, las de la Etapa 2 del Split se citan como (E2-1) a (E2-7).

\startitemize
\item {\bf Monolítico:} decide a la vez las rutas ($z_{ijb}$, $q_{ijb}$ y el tiempo de ruta
$r_b$, restricciones (1)--(13)), la asignación $x_{pb}$, la duración $u_{pb}$ de cada lote en
cada picker ((14)--(17)), las terminaciones $c_{pb}$ y $c_b$ ((18)), la secuencia $y_{bk}$
((19)--(20)) y el atraso y adelanto de cada orden ((21)--(22)). La cadena de frío es una cota de
la variable: $u_{pb} \le T^{cold} + \varepsilon$.
\item {\bf Split, Etapa 1:} las mismas restricciones (1)--(13), resueltas por lote, minimizando
$\sum_b r_b$. Su resultado es el tiempo de ruta $R_b$ de cada lote.
\item {\bf Split, Etapa 2:} con $R_b$ como dato, la duración de cada par es un parámetro,
$T^{exec}_{pb} = F_b + R_b / V_p$, y solo existen los pares elegibles por la cadena de frío,
$\mathcal{PB}^{elig} = \{(p,b) : T^{exec}_{pb} \le T^{cold} + \varepsilon\}$. Restricciones:
asignación (E2-1), $c_{pb} \ge T^{exec}_{pb} \, x_{pb}$ (E2-2), $c_b$ (E2-3), secuencia (E2-4) y
(E2-5), atraso y adelanto (E2-6) y (E2-7).
\item {\bf Mismo objetivo:} $W^{c} \sum_b c_b + W^{pen} \left( W^{e} \sum_n ea_n + W^{t} \sum_n
ta_n \right)$ en el monolítico y en la Etapa 2. Ninguno de los dos tiene un término de ruta.
\stopitemize
\stopsubsection

\startsubsection[title={Correspondencia de restricciones}]

\blank[medium]
\bTABLE[split=yes, align=middle, frame=on, width=broad]
\setupTABLE[row][1][background=color, backgroundcolor=darkgray, style=bold]
\setupTABLE[column][1][align=flushleft]
\setupTABLE[column][2][align=flushleft]
\bTR \bTH Monolítico (nombre en el código) \eTH \bTH Split \eTH \eTR
\bTR \bTD (1)--(13) ruteo y flujo, de \ruta{start_from_depot} a \ruta{define_route_time} \eTD \bTD Etapa 1: las mismas restricciones, una copia por lote \eTD \eTR
\bTR \bTD (14) \ruta{assign_exactly_one_picker}, sobre todos los pickers \eTD \bTD (E2-1), solo sobre los pickers elegibles \eTD \eTR
\bTR \bTD (15) \ruta{batch_execution_time_lower} \eTD \bTD Absorbida en (E2-2), con $T^{exec}_{pb}$ como dato \eTD \eTR
\bTR \bTD (16) \ruta{batch_execution_time_upper} \eTD \bTD No existe: la Etapa 2 no tiene $u_{pb}$ \eTD \eTR
\bTR \bTD (17) \ruta{processing_le_completion} \eTD \bTD Absorbida en (E2-2) \eTD \eTR
\bTR \bTD Cota $u_{pb} \le T^{cold} + \varepsilon$ \eTD \bTD Dominio $\mathcal{PB}^{elig}$, más un pre-chequeo en Python si un lote queda sin pickers elegibles \eTD \eTR
\bTR \bTD (18) \ruta{c_b_lower_bound} \eTD \bTD (E2-3), solo pares elegibles \eTD \eTR
\bTR \bTD (19)--(20) \ruta{non_overlap_1} y \ruta{non_overlap_2} \eTD \bTD (E2-4)--(E2-5), con $T^{exec}$ como dato y solo si los dos pares son elegibles \eTD \eTR
\bTR \bTD (21)--(22) \ruta{tardiness_definition} y \ruta{earliness_definition} \eTD \bTD (E2-6)--(E2-7), solo pares elegibles \eTD \eTR
\eTABLE
\blank[medium]

Los valores de $M$ y $M^{seq}$ son los mismos en los dos modelos: los dos los leen del mismo
archivo de parámetros que escribe el constructor de datos. También coinciden los nombres de las
variables ($z_{ijb}$, $x_{pb}$, $y_{bk}$, $c_{pb}$, $c_b$, $r_b$), lo que permite comparar
soluciones variable por variable.
\stopsubsection

\startsubsection[title={La Etapa 1 se separa por lote}]

Las restricciones (1)--(13) son \quote{para cada lote}: ninguna vincula variables de dos lotes
distintos, y el objetivo de la Etapa 1 es una suma de un término por lote. Por lo tanto,
minimizar $\sum_b r_b$ es lo mismo que minimizar cada $r_b$ por separado, y la Etapa 1 da, para
cada lote, su ruta mínima $R^{*}_b$. Es lo que hace la implementación al resolver un MIP por lote.
Además, ninguna restricción de ruteo depende del picker ni del orden entre lotes, así que la ruta
mínima de un lote no depende de quién lo ejecute.
\stopsubsection

\startsubsection[title={Demostración}, reference=sec:demostracion]

Sean $z^{M}$ el valor óptimo del monolítico y $z^{S}$ el del Split (la Etapa 2 resuelta con las
rutas mínimas de la Etapa 1).

{\bf Sentido 1: de una solución del Split a una del monolítico ($z^{M} \le z^{S}$).} Se toma una
solución factible del Split: rutas de la Etapa 1 con tiempos $R^{*}_b$, y $(x, c, c_b, y, ta,
ea)$ de la Etapa 2. Se arma una solución del monolítico con las mismas rutas ($r_b = R^{*}_b$),
los mismos valores de $x$, $y$, $c$, $c_b$, $ta$ y $ea$ en los pares elegibles, $x_{pb} = 0$ y
$c_{pb} = 0$ en los no elegibles, y $u_{pb} = T^{exec}_{pb}$ si $x_{pb} = 1$, o $0$ si no.

\startitemize
\item {\bf (1)--(13):} son las mismas restricciones que la Etapa 1.
\item {\bf (14):} cada lote tiene exactamente un picker elegible asignado (E2-1), y los no elegibles
valen 0.
\item {\bf (15)--(16):} si $x_{pb} = 1$, $u_{pb} = F_b + r_b / V_p$ las cumple con igualdad; si
$x_{pb} = 0$, $u_{pb} = 0$ las cumple siempre que $M \ge F_b + r_b / V_p$.
\item {\bf Cadena de frío:} solo se asignan pares elegibles, así que $u_{pb} = T^{exec}_{pb} \le
T^{cold} + \varepsilon$.
\item {\bf (17):} $u_{pb} = T^{exec}_{pb} \, x_{pb} \le c_{pb}$ es exactamente (E2-2).
\item {\bf (18)--(22):} con $x_{pb} = 1$ coinciden con (E2-3)--(E2-7), porque $F_k + r_k / V_p =
T^{exec}_{pk}$; con $x_{pb} = 0$ quedan desactivadas por la Big-M.
\item {\bf Objetivo:} la misma expresión sobre los mismos valores de $c_b$, $ea$ y $ta$.
\stopitemize

{\bf Sentido 2: de una solución del monolítico a una del Split ($z^{S} \le z^{M}$).} Se toma una
solución óptima del monolítico y se reemplaza la ruta de cada lote por la de la Etapa 1. Como la
ruta del monolítico cumple (1)--(13), es factible para la Etapa 1 de ese lote, y por lo tanto
$R^{*}_b \le r_b$. Se conservan $x$, $y$, $c$, $c_b$, $ta$ y $ea$.

\startitemize
\item {\bf Elegibilidad:} si $x_{pb} = 1$, (15)--(16) fijan $u_{pb} = F_b + r_b / V_p \le
T^{cold} + \varepsilon$. Como $R^{*}_b \le r_b$, también $T^{exec}_{pb} = F_b + R^{*}_b / V_p \le
T^{cold} + \varepsilon$: el par es elegible, así que la asignación existe en la Etapa 2 y cumple
(E2-1).
\item {\bf (E2-2):} $c_{pb} \ge u_{pb} = F_b + r_b / V_p \ge F_b + R^{*}_b / V_p =
T^{exec}_{pb}$.
\item {\bf (E2-4)--(E2-5):} la separación que exigen entre dos lotes del mismo picker es la
duración del segundo, que con la ruta mínima es igual o menor. Si el monolítico la cumplía, la
Etapa 2 también.
\item {\bf (E2-3), (E2-6) y (E2-7):} son las mismas que (18), (21) y (22), sobre los mismos valores.
\item {\bf Objetivo:} no cambia, porque no depende de las rutas.
\stopitemize

{\bf Conclusión.} $z^{M} = z^{S}$. Las mismas dos construcciones muestran que los dos modelos
coinciden en factibilidad: si el Split tiene una solución, el monolítico también (sentido 1); si el
monolítico tiene una, la Etapa 1 es factible para cada lote y la Etapa 2 también (sentido 2).

{\bf Un matiz sobre la frase original.} Se decía que una ruta más corta \quote{adelanta las
terminaciones}. No es necesariamente así: con espera voluntaria, al acortar la ruta las
terminaciones no se mueven. Lo que hace la ruta más corta es dejar factibles las mismas
terminaciones, y eso es lo que usa la demostración.
\stopsubsection

\startsubsection[title={Condiciones}, reference=sec:condiciones]

La demostración vale en aritmética exacta y bajo cuatro condiciones:

\startitemize[n]
\item {\bf $M$ suficientemente grande:} cuando $x_{pb} = 0$, las restricciones con $M$ ((15), (16),
(18), (21) y (22), y sus equivalentes de la Etapa 2) tienen que quedar desactivadas. En
particular, $M \ge F_b + r_b / V_p$. Los dos modelos usan el mismo $M$, que es un dato del
escenario: entre $10^{3}$ y $10^{5}$ en las instancias funcionales y $10^{6}$ en el perfil
realista de stress.
\item {\bf $M^{seq}$ suficientemente grande:} cubre la diferencia entre dos terminaciones del mismo
picker. Se calcula por instancia como una cota del horizonte, con un piso de 500, y es el mismo
valor en los dos modelos. Con 500 fijo, el Split llegó a declarar infactible una instancia de 40
lotes en un picker (ver las notas de desarrollo). La cota cubre la ejecución de los lotes, pero no
la espera voluntaria hasta las fechas de entrega; como los dos modelos usan el mismo valor, eso
no afecta la equivalencia entre ellos.
\item {\bf Etapa 1 resuelta a optimalidad:} si la Etapa 1 se cortara por tiempo, $R_b$ podría no
ser la ruta mínima y el Split podría dar un objetivo mayor que el monolítico. En los tests la
Etapa 1 corre sin límite de tiempo. Solo el experimento a igual tiempo le pone un presupuesto, y
ahí la diferencia es buscada.
\item {\bf Aritmética exacta:} en la práctica, SCIP trabaja con tolerancias de factibilidad e
integralidad del orden de $10^{-6}$. La consecuencia más concreta está en la cadena de frío: en el
Split es un filtro exacto sobre los datos ($T^{exec}_{pb} \le T^{cold} + 10^{-9}$), y en el
monolítico es una cota de una variable que SCIP respeta con su tolerancia. Muy cerca del borde
(entre $10^{-9}$ y unos $10^{-6}$ por encima del límite), el monolítico podría aceptar un par que
el Split rechaza. Está documentado y no tiene test (sección \in[sec:limitaciones]).
\stopitemize
\stopsubsection

\startsubsection[title={Por qué las soluciones óptimas pueden ser distintas}, reference=sec:porque_alternativos]

La demostración iguala los valores óptimos, no los conjuntos de soluciones óptimas. Hay tres
fuentes de diferencia:

\startitemize
\item {\bf Rutas no mínimas:} el sentido 2 muestra que una solución óptima del monolítico con rutas
más largas que las mínimas tiene una contraparte del Split con el mismo objetivo. El Split siempre
devuelve rutas mínimas, porque la Etapa 1 las minimiza; el monolítico puede devolver cualquier
ruta que no empeore el objetivo.
\item {\bf Espera voluntaria:} si un lote igual va a esperar hasta su fecha de entrega, una ruta más
larga que llega a tiempo cuesta lo mismo. Es la razón por la que el monolítico elige rutas más
largas en las instancias de la sección \in[sec:alternativos].
\item {\bf Empates:} cuando dos asignaciones de lotes a pickers, o dos rutas del mismo largo, dan el
mismo objetivo, cada resolución rompe el empate a su manera.
\stopitemize
\stopsubsection

\stopsection

%%%%%%%%%%%%%%%%%%%%%%%%%%%%%%%%%%%%%%%%%%%%%%%%%%%%%%%%%%%%%%%%%%%%
\startsection[title={Línea 2: la suite de tests compartida}, reference=sec:tests]

\startsubsection[title={Misma instancia en los dos modelos}]

La suite funcional tiene 38 archivos en cada modelo, con los mismos números de test. Hay dos
clases de archivos, y en las dos la instancia sale de la misma llamada:

\startitemize
\item {\bf Heredados (30 archivos):} son los tests de la versión anterior del modelo, que no tenía
velocidades de picker ni cadena de frío. Cada archivo de v5 es una línea que envuelve el mismo
módulo de test en los dos modelos. Un adaptador reemplaza la resolución original por la del
modelo v5, con $V_p = 1$ y la cadena de frío desactivada ($T^{cold} = 10^{9}$), y conserva el
escenario y las aserciones sin tocarlas. El adaptador del Split reemplaza además la aserción de
no solapamiento, que en el original se calcula con $u_{pb}$, una variable que la Etapa 2 no tiene;
el del monolítico solo cambia el archivo \ruta{.zpl}.
\item {\bf Propios de v5 (archivos 06, 23, 24 y 34 a 38):} los escenarios están en un módulo común
y los dos modelos llaman al mismo constructor con los mismos argumentos. Los archivos 06, 35, 36 y
37 son idénticos salvo la clase base; el 23 difiere en cómo se desempaqueta el resultado, el 34
solo en comentarios, y el 38 llega al resultado por otro camino (el Split corre solo la Etapa 1,
que ya es infactible). El 24b existe solo en el Split (ver más abajo).
\stopitemize

En stress, los 19 tests de cada modelo llaman a la misma función con los mismos argumentos; se
verificó textualmente para los 19 pares.

Las dos clases base exponen los mismos ayudantes de aserción y devuelven el resultado con la misma
forma, de modo que el cuerpo de un test propio es el mismo en los dos modelos.
\stopsubsection

\startsubsection[title={Qué afirman los tests}]

\startitemize
\item {\bf Stress:} estado óptimo o factible, rutas que salen del depósito, llegan al final y
cubren las ubicaciones requeridas, cada lote asignado y sin solapamiento entre lotes del mismo
picker (tolerancia $10^{-5}$). No afirman el objetivo.
\item {\bf Heredados:} estado, y en varios casos el grafo exacto de la ruta. Ninguno afirma el
objetivo. Los que exigen una ruta puntual pueden fallar en un modelo correcto si hay rutas
alternativas igual de óptimas (sección \in[sec:test06]).
\item {\bf Propios de v5:} valores exactos calculados a mano sobre la geometría del escenario
(tabla siguiente).
\stopitemize

\blank[medium]
\bTABLE[split=yes, align=middle, frame=on, width=broad]
\setupTABLE[row][1][background=color, backgroundcolor=darkgray, style=bold]
\setupTABLE[column][1][width=1.4cm]
\setupTABLE[column][2][align=flushleft]
\setupTABLE[column][3][width=2.6cm]
\bTR \bTH Test \eTH \bTH Qué afirma \eTH \bTH ¿Objetivo exacto? \eTH \eTR
\bTR \bTD 06 \eTD \bTD 06b: con fecha 10, cada lote toma la ruta corta, termina en 10, objetivo 20. 06a: con fecha 15, cobertura, cada lote termina en 15, objetivo 30; la ruta no se afirma. \eTD \bTD Sí \eTD \eTR
\bTR \bTD 23 \eTD \bTD Tiempos de ruta 5 y 14; todos los lotes asignados y sin solapamiento; el lote de ruta larga va al picker rápido. \eTD \bTD No \eTD \eTR
\bTR \bTD 24 \eTD \bTD 24a: la cadena de frío deja solo al picker rápido. 24b (solo Split): tras la Etapa 1, ningún par elegible. \eTD \bTD No \eTD \eTR
\bTR \bTD 34 \eTD \bTD Tiempos de ruta 5 y 14; asignación; terminaciones 7 y 9 por la fórmula; un par justo en el límite de frío sigue elegible; sin pickers elegibles, infactible. \eTD \bTD No \eTD \eTR
\bTR \bTD 35 \eTD \bTD Frío por lote y no por turno (objetivo 27); tolerancia de $5 \cdot 10^{-10}$; la parte fija cuenta para el frío (infactible). \eTD \bTD Sí \eTD \eTR
\bTR \bTD 36 \eTD \bTD Velocidad 4 (objetivo 16,5), velocidad 2 (25), velocidad 0,5 (atraso 10, objetivo 310); velocidad 0 o negativa se rechaza nombrando al picker. \eTD \bTD Sí \eTD \eTR
\bTR \bTD 37 \eTD \bTD Con $W^{c} = 1$ el lote espera hasta la fecha (objetivo 20); con $W^{c} = 5$ no espera (objetivo 78). \eTD \bTD Sí \eTD \eTR
\bTR \bTD 38 \eTD \bTD Una ruta que tendría que pasar dos veces por un vértice es infactible. \eTD \bTD -- \eTD \eTR
\eTABLE
\blank[medium]

Donde hay varias terminaciones, los tests las comparan ordenadas y no atadas a un lote puntual, y
el 36b solo verifica qué pickers se usan: así toleran empates que no cambian el objetivo. La
igualdad de objetivos entre los dos modelos no se verifica dentro de los tests sino después, con
la comparación cruzada (sección \in[sec:cruzada]).
\stopsubsection

\startsubsection[title={Cantidad de tests}]

\blank[medium]
\bTABLE[split=yes, align=middle, frame=on, width=fit]
\setupTABLE[row][1][background=color, backgroundcolor=darkgray, style=bold]
\bTR \bTH \eTH \bTH Split \eTH \bTH Monolítico \eTH \eTR
\bTR \bTD Métodos heredados \eTD \bTD 45 \eTD \bTD 45 \eTD \eTR
\bTR \bTD Métodos propios de v5 \eTD \bTD 19 \eTD \bTD 18 \eTD \eTR
\bTR \bTD Total funcional \eTD \bTD 64 \eTD \bTD 63 \eTD \eTR
\bTR \bTD Stress \eTD \bTD 19 \eTD \bTD 19 \eTD \eTR
\eTABLE
\blank[medium]

El único test sin equivalente es el 24b: verifica el pre-chequeo de la cadena de frío del Split
después de correr solo la Etapa 1, y el monolítico no tiene etapas.
\stopsubsection

\startsubsection[title={El test heredado 06: un fallo que no era un error}, reference=sec:test06]

El test 06 heredado exigía la ruta más corta (L1→L4→L5, tiempo 4) en una instancia con fecha de
entrega 15. En v5, con espera voluntaria, la ruta L1→L2→L4→L5 (tiempo 7,5) también llega a
tiempo, y los dos lotes se declaran terminados en 15 de todos modos: las dos rutas dan objetivo
30. El monolítico elegía la larga y el test fallaba; el Split pasaba solo porque su Etapa 1
devuelve la ruta mínima por construcción. Hasta el 29-09 figuró como \quote{falla conocida} del
monolítico.

El 29-09 se reemplazó por un test propio de v5 sobre la misma instancia, con dos casos: con fecha
10 la ruta larga atrasa, así que la corta es la única óptima, y se afirma (06b); con fecha 15
toda ruta llega, así que se afirman el objetivo y las terminaciones, no la ruta (06a). Los dos
modelos pasan los dos casos, y la comparación cruzada marca el 06a como óptimo alternativo. Es
el ejemplo más claro de por qué la validación separa el objetivo de la solución.
\stopsubsection

\startsubsection[title={Resultado}]

Con el código actual, el Split pasa 63 de sus 64 tests funcionales y el monolítico 62 de 63. El
único que no pasa es el mismo en los dos: el 30a, con capacidades fraccionarias y 10 de sus 12
lotes vacíos, da timeout a los 25 minutos en los dos modelos. Al verificar una solución de esa
instancia se confirmó la causa: con los lotes vacíos libres, SCIP no cierra en 600~s; sin ellos,
cierra en 0,13~s. Son $10!$ permutaciones indistinguibles de lotes vacíos.
\stopsubsection

\stopsection

%%%%%%%%%%%%%%%%%%%%%%%%%%%%%%%%%%%%%%%%%%%%%%%%%%%%%%%%%%%%%%%%%%%%
\startsection[title={Línea 3: comparación cruzada de soluciones}, reference=sec:cruzada]

\startsubsection[title={Cómo se guardan las soluciones}]

Con la variable de entorno \ruta{WRITE_SOLUTIONS=1}, el runner de tests guarda cada solución en un
JSON, después de medir la duración del test (no afecta los tiempos). El CSV de resultados enlaza
cada fila con su JSON (columna \ruta{solution_file}). Cada JSON tiene:

\startitemize
\item {\bf La instancia completa:} todos los conjuntos y parámetros que se escribieron en los
\ruta{.dat}, tal como los recibió el modelo.
\item {\bf La solución:} estado, objetivo y todos los valores no nulos de las variables, con su
nombre de ZIMPL (por ejemplo, \ruta{x_pb$P1$B2}).
\item {\bf En el Split, las dos etapas:} los valores de la Etapa 2 y, aparte, los de la Etapa 1
($z_{ijb}$, $q_{ijb}$, $r_b$).
\stopitemize

Las corridas del 29-09 dejaron 119 soluciones: 54 funcionales del monolítico, 56 funcionales del
Split y 9 de stress del Split. El monolítico no tiene soluciones de stress porque no resuelve
ninguno.
\stopsubsection

\startsubsection[title={El comparador}]

\ruta{verification/compare_v5.py} lee los CSV finales de los dos modelos y los JSON enlazados, y no
resuelve nada. Para cada test:

\startitemize[n]
\item {\bf Empareja} la fila del Split con la del monolítico por el nombre del método de test.
\item {\bf Compara el estado.} Si alguno no terminó, el veredicto es \quote{mismo estado} (los dos
con timeout) o \quote{no resueltos los dos}.
\item {\bf Compara el objetivo} con la tolerancia de la sección \in[sec:objetivo]. Si los dos
están vacíos, es \quote{infactible en los dos}. Si el Split se detuvo antes de la Etapa 2 (Etapa 1
infactible), su objetivo no se usa, porque sería un tiempo de ruta.
\item {\bf Si el objetivo coincide, compara la solución,} tomando como activo todo binario mayor que
0,5:
\startitemize
\item {\bf Rutas:} el conjunto de arcos activos $z_{ijb}$ de cada lote.
\item {\bf Asignación:} el conjunto de pares activos $x_{pb}$.
\item {\bf Orden:} si la asignación coincide, la lista de lotes de cada picker ordenada por
$c_{pb}$.
\stopitemize
\item {\bf Clasifica:} \quote{misma solución} si todo coincide; \quote{óptimo alternativo} si el
objetivo coincide y algo de lo anterior no; \quote{objetivo distinto} si no coincide el objetivo.
\stopitemize

El comparador agrupa los tests en heredados, propios de v5 y stress, escribe un CSV con el
veredicto y el detalle de cada fila, y termina con error solo si algún objetivo difiere: es la
condición que rompería la equivalencia.
\stopsubsection

\startsubsection[title={Resultados}, reference=sec:res_cruzada]

\blank[medium]
\bTABLE[split=yes, align=middle, frame=on, width=broad]
\setupTABLE[row][1][background=color, backgroundcolor=darkgray, style=bold]
\bTR \bTH Grupo \eTH \bTH Filas \eTH \bTH Misma solución \eTH \bTH Óptimo alternativo \eTH \bTH Infactible en los dos \eTH \bTH Otro \eTH \bTH Objetivo distinto \eTH \eTR
\bTR \bTD Heredados \eTD \bTD 44 \eTD \bTD 15 \eTD \bTD 20 \eTD \bTD 8 \eTD \bTD 1 (30a, timeout en los dos) \eTD \bTD 0 \eTD \eTR
\bTR \bTD Propios de v5 \eTD \bTD 19 \eTD \bTD 13 \eTD \bTD 1 (06a) \eTD \bTD 4 \eTD \bTD 1 (24b, solo Split) \eTD \bTD 0 \eTD \eTR
\bTR \bTD {\bf Funcional} \eTD \bTD {\bf 63} \eTD \bTD {\bf 28} \eTD \bTD {\bf 21} \eTD \bTD {\bf 12} \eTD \bTD {\bf 2} \eTD \bTD {\bf 0} \eTD \eTR
\bTR \bTD Stress \eTD \bTD 19 \eTD \bTD -- \eTD \bTD -- \eTD \bTD -- \eTD \bTD 9 timeout en los dos; 10 resueltos solo por el Split o sin resultado \eTD \bTD -- \eTD \eTR
\eTABLE
\blank[medium]

\startitemize
\item {\bf Ningún objetivo distinto.} Los 49 pares resueltos por los dos tienen el mismo objetivo.
La mayor diferencia es de $3{,}83 \cdot 10^{-6}$ en valor absoluto ($1{,}2 \cdot 10^{-9}$
relativo), en los tests 13 y 14: 3.086,9999999999754 en el Split contra 3.086,9999961653775 en el
monolítico. La siguiente es de $6{,}1 \cdot 10^{-10}$ (25b).
\item {\bf Mismo estado en todos los infactibles.} Los 12 infactibles lo son en los dos modelos. En
5 de ellos (07, 21, 26a, 31a y 38) el Split ya se detiene en la Etapa 1, porque alguna ruta es
imposible, y el monolítico llega a la misma conclusión resolviendo todo junto. El 36c cuenta acá
aunque no es un infactible propiamente dicho: los dos modelos rechazan la entrada (velocidad no
positiva) antes de resolver.
\item {\bf Stress sin pares comparables.} El monolítico no termina ninguno de los 19. El Split
termina 9, y para esos la comparación de objetivos no es posible: es el hueco que cubre la línea 4.
\stopitemize

{\bf La base comparable.} El documento de comparación del repositorio de implementación usa la
misma base que la comparación con las heurísticas: 49 pares resueltos y 11 infactibles, y deja
fuera el 24b (solo Split), el 30a (timeout, sin referencia), el 36c (entrada rechazada) y el 37a
(verifica la espera voluntaria, que las heurísticas no modelan; entre los dos modelos exactos sí
es comparable y da el mismo objetivo, 20). Por eso cuenta 11 infactibles y no 12. Que también dé 49
resueltos es una coincidencia de dos bases distintas: suma el 02 y resta el 37a.
\stopsubsection

\startsubsection[title={Los 21 óptimos alternativos}, reference=sec:alternativos]

\blank[medium]
\bTABLE[split=yes, align=middle, frame=on, width=broad]
\setupTABLE[row][1][background=color, backgroundcolor=darkgray, style=bold]
\setupTABLE[column][2][align=flushleft]
\bTR \bTH Caso \eTH \bTH Test \eTH \bTH Difiere en \eTH \bTH $\sum_b r_b$ Split \eTH \bTH $\sum_b r_b$ mono \eTH \eTR
\bTR \bTD 06a \eTD \bTD Todas las rutas llegan a tiempo \eTD \bTD rutas \eTD \bTD 8 \eTD \bTD 15 \eTD \eTR
\bTR \bTD 15 \eTD \bTD Tres órdenes, un picker \eTD \bTD rutas \eTD \bTD 20 \eTD \bTD 20 \eTD \eTR
\bTR \bTD 19a \eTD \bTD Ruta extendida, ubicación dup. \eTD \bTD rutas \eTD \bTD 7 \eTD \bTD 10 \eTD \eTR
\bTR \bTD 19b \eTD \bTD Ruta extendida, penalizaciones \eTD \bTD rutas \eTD \bTD 7 \eTD \bTD 10 \eTD \eTR
\bTR \bTD 19c \eTD \bTD Ruta extendida \eTD \bTD rutas \eTD \bTD 7 \eTD \bTD 10 \eTD \eTR
\bTR \bTD 13 \eTD \bTD Dos períodos, grilla \eTD \bTD rutas y asignación \eTD \bTD 72 \eTD \bTD 72 \eTD \eTR
\bTR \bTD 14 \eTD \bTD Un período, grilla \eTD \bTD rutas y asignación \eTD \bTD 72 \eTD \bTD 72 \eTD \eTR
\bTR \bTD 16 \eTD \bTD Seis órdenes, dos pickers \eTD \bTD rutas y asignación \eTD \bTD 44 \eTD \bTD 44 \eTD \eTR
\bTR \bTD 25b \eTD \bTD Grilla: ubicación dup. \eTD \bTD rutas y asignación \eTD \bTD 34 \eTD \bTD 34 \eTD \eTR
\bTR \bTD 25e \eTD \bTD Grilla: penalización por fecha \eTD \bTD rutas y asignación \eTD \bTD 36 \eTD \bTD 45 \eTD \eTR
\bTR \bTD 25f \eTD \bTD Grilla: orden igual al carro \eTD \bTD rutas y asignación \eTD \bTD 36 \eTD \bTD 36 \eTD \eTR
\bTR \bTD 26d \eTD \bTD Puente: atraso \eTD \bTD rutas y asignación \eTD \bTD 34 \eTD \bTD 34 \eTD \eTR
\bTR \bTD 29 \eTD \bTD Lote vacío permitido \eTD \bTD rutas y asignación \eTD \bTD 52 \eTD \bTD 100 \eTD \eTR
\bTR \bTD 11 \eTD \bTD Cuatro órdenes \eTD \bTD asignación \eTD \bTD \eTD \bTD \eTD \eTR
\bTR \bTD 20 \eTD \bTD Grilla multiubicación \eTD \bTD asignación \eTD \bTD \eTD \bTD \eTD \eTR
\bTR \bTD 25c \eTD \bTD Grilla multiubicación con capacidades \eTD \bTD asignación \eTD \bTD \eTD \bTD \eTD \eTR
\bTR \bTD 25g \eTD \bTD Grilla: ubic. compartidas \eTD \bTD asignación \eTD \bTD \eTD \bTD \eTD \eTR
\bTR \bTD 26b \eTD \bTD Puente: adelanto \eTD \bTD asignación \eTD \bTD \eTD \bTD \eTD \eTR
\bTR \bTD 26c \eTD \bTD Puente: ubic. compartidas \eTD \bTD asignación \eTD \bTD \eTD \bTD \eTD \eTR
\bTR \bTD 27 \eTD \bTD Pesos distintos \eTD \bTD asignación \eTD \bTD \eTD \bTD \eTD \eTR
\bTR \bTD 33c \eTD \bTD Servicio cero \eTD \bTD asignación \eTD \bTD \eTD \bTD \eTD \eTR
\eTABLE
\blank[medium]

Las sumas de tiempos de ruta salen de los valores de $r_b$ de las soluciones guardadas. Se
calcularon al escribir estas notas para comprobar lo que predice la sección
\in[sec:porque_alternativos]:

\startitemize
\item {\bf El monolítico nunca encontró una ruta más corta que el Split.} Es la comprobación empírica
de que la Etapa 1 da la ruta mínima.
\item {\bf En 7 de los 13 pares con rutas distintas, el tiempo total de ruta es el mismo:} son rutas
distintas del mismo largo, es decir, empates.
\item {\bf En los otros 6, el monolítico eligió rutas más largas} con el mismo objetivo (06a, 19a,
19b, 19c, 25e y 29). Es la consecuencia directa de la espera voluntaria y de que el objetivo no
tenga un término de ruta.
\item {\bf En los 8 que difieren solo en la asignación,} las rutas son idénticas y la diferencia es
qué picker toma cada lote: empates de asignación.
\stopitemize

En los 28 pares con la misma solución, coinciden rutas, asignación y orden de cada picker.
\stopsubsection

\stopsection

%%%%%%%%%%%%%%%%%%%%%%%%%%%%%%%%%%%%%%%%%%%%%%%%%%%%%%%%%%%%%%%%%%%%
\startsection[title={Línea 4: el monolítico como verificador}, reference=sec:verificador]

\startsubsection[title={Para qué}]

En stress el monolítico no termina ninguna instancia en 40 minutos, así que no hay un óptimo del
monolítico contra el cual comparar. Se cambia la pregunta: la solución que devuelve el Split, ¿es
factible bajo las restricciones del monolítico, y le cuesta lo mismo? Si la Etapa 1 y la Etapa 2
implementan bien las restricciones del monolítico, la respuesta tiene que ser sí. Lo que esta línea
detecta son errores de implementación en las etapas del Split: una restricción mal escrita, un dato
mal pasado de una etapa a la otra o una cadena de frío mal filtrada. El mismo verificador se usó
para validar las soluciones de las heurísticas.
\stopsubsection

\startsubsection[title={Cómo funciona}]

\ruta{verification/mono_compliance.py} toma el \ruta{.zpl} del monolítico sin modificar, escribe
los datos de la instancia, lo compila y fija sobre las variables del modelo compilado los valores
de la solución, buscándolas por su nombre de ZIMPL. No copia ni reescribe ninguna restricción: el
\ruta{.zpl} del monolítico es la única definición de qué es una solución válida. Si un nombre no
existe en el modelo compilado, el verificador falla en vez de ignorarlo. La comprobación tiene
tres capas:

\startitemize
\item {\bf Capa 0, chequeos en Python:} cada lote asignado a exactamente un picker; rutas que solo
usan nodos y arcos del grafo, sin revisitas; ubicaciones requeridas cubiertas; cadena de frío
($T^{exec} \le T^{cold} + 10^{-9}$); tiempos de ejecución, terminaciones y objetivo reportados
iguales a los que da la fórmula del monolítico recalculada en Python.
\item {\bf Nivel 1, las decisiones:} se fijan las rutas ($z_{ijb}$ de todos los arcos), la
asignación ($x_{pb}$ de todos los pares) y la secuencia ($y_{bk}$), y se dejan libres los tiempos.
SCIP resuelve: si es infactible, las decisiones no son válidas para el monolítico.
\item {\bf Nivel 2, la solución exacta:} se fijan además las terminaciones $c_{pb}$ y $c_b$
reportadas. SCIP comprueba que la solución entera sea factible y calcula su objetivo.
\stopitemize

El veredicto es {\bf \quote{compliant}} solo si no hay hallazgos en la capa 0, los dos niveles son
factibles y el objetivo del nivel 2 coincide con el reportado (tolerancia relativa $10^{-6}$).
Cualquier otro resultado es \quote{violations}, con el detalle. Si el nivel 1 es infactible, el
verificador libera por turnos las rutas, la secuencia o la asignación, y con hasta 40 lotes
también la ruta de cada lote por separado, para señalar qué decisión es la inválida.

Algunos detalles de diseño:

\startitemize
\item {\bf Big-M:} si el $M$ o el $M^{seq}$ del escenario fueran demasiado chicos para la solución
que se verifica, se reemplazan por un valor válido y queda anotado; nunca cuenta como violación.
\item {\bf Lotes vacíos:} el monolítico obliga a rutear, asignar y cobrar todo lote de
$\mathcal{B}$, aunque no tenga órdenes, mientras que las heurísticas los ignoran. Por decisión del
equipo, el verificador saca de $\mathcal{B}$ los lotes sin órdenes que la solución no usa, y lo
anota. Ninguna restricción cambia.
\item {\bf Instancia de 100×100:} compilar el monolítico completo llevó a ZIMPL a unos 5~GB sin
terminar. Cuando todas las rutas están fijas, el verificador trabaja sobre el subgrafo de esas
rutas (sus arcos, sus nodos, el depósito y el final), que es exactamente equivalente con las rutas
fijas. El control: las 248 soluciones que ya estaban verificadas dieron el mismo veredicto, los
mismos estados y los mismos objetivos con y sin el recorte.
\item {\bf Adaptación de una solución del Split:} las rutas se reconstruyen a partir de los arcos
activos de la Etapa 1, el orden de cada picker a partir de $c_{pb}$ y las terminaciones de los
lotes a partir de $c_b$.
\stopitemize
\stopsubsection

\startsubsection[title={Prueba de que el verificador puede fallar}]

Un verificador que siempre dice \quote{compliant} no prueba nada. Sobre una solución chica se
introdujeron cuatro errores a propósito: un arco que no existe en el grafo, dos terminaciones
superpuestas del mismo picker, una violación de la cadena de frío y un objetivo mal reportado. El
verificador detectó los cuatro.
\stopsubsection

\startsubsection[title={Resultados}]

\startitemize
\item {\bf Split en stress:} las 9 soluciones de stress del Split dieron \quote{compliant} (corrida
del 29-09 con \ruta{python -m verification.compare_v5 --verify-split-stress}). Ese veredicto quedó
en el registro de cambios del repositorio, pero no en el CSV de comparación guardado, que se
generó sin esa opción. Se reproduce con el mismo comando.
\item {\bf Heurísticas, como contexto:} las 150 soluciones de HC y ACO (56 funcionales y 19 de stress
por heurística) también dieron \quote{compliant}. En la primera verificación, sobre una suite
anterior, el nivel 1 bajó el objetivo de 7 de 51 casos entre 7,9~\% y 93,6~\% con las mismas
decisiones: el monolítico espera hasta la fecha de entrega y las heurísticas no. No es una
violación; es la espera voluntaria medida desde otro lado.
\stopitemize
\stopsubsection

\stopsection

%%%%%%%%%%%%%%%%%%%%%%%%%%%%%%%%%%%%%%%%%%%%%%%%%%%%%%%%%%%%%%%%%%%%
\startsection[title={Línea 5: pares en la exploración}, reference=sec:pares]

\startsubsection[title={Método}]

Los estudios 08 a 11 de la exploración del monolítico usaron el {\em modo par}: cada instancia se
genera una vez y se resuelve primero con el Split y después con el monolítico, con el mismo tope.
Los estudios 08 y 11 usan el perfil realista (pesos 4/5, velocidades en bandas, 26 órdenes por
lote); el 09 repite instancias de los estudios 01 a 03 del Split, y el 10 varía un factor por vez
(órdenes, pesos, velocidades, lotes, pickers, grilla). Son más grandes que las funcionales: hasta
6 lotes y grillas de hasta 7×7, y un punto de 10×10. El objetivo de cada modelo se registra con dos
decimales.

\blank[medium]
\bTABLE[split=yes, align=middle, frame=on, width=broad]
\setupTABLE[row][1][background=color, backgroundcolor=darkgray, style=bold]
\setupTABLE[column][2][align=flushleft]
\bTR \bTH Estudio \eTH \bTH Qué varía \eTH \bTH Tope por modelo \eTH \bTH Pares \eTH \bTH Óptimo probado por los dos \eTH \bTH Mismo objetivo \eTH \eTR
\bTR \bTD 08 \eTD \bTD Frontera: lotes {1, 2, 3, 4, 6} × lado {3 a 7} \eTD \bTD 600~s \eTD \bTD 25 \eTD \bTD 15 \eTD \bTD 14 \eTD \eTR
\bTR \bTD 09 \eTD \bTD Las instancias de los estudios 01 a 03 \eTD \bTD 330~s \eTD \bTD 15 \eTD \bTD 0 \eTD \bTD -- \eTD \eTR
\bTR \bTD 10 \eTD \bTD Un factor por vez desde 4 lotes, 2 pickers, 5×5 \eTD \bTD 630~s \eTD \bTD 27 \eTD \bTD 25 \eTD \bTD 25 \eTD \eTR
\bTR \bTD 11 \eTD \bTD Tope de 40 minutos en los puntos difíciles \eTD \bTD 2.430~s \eTD \bTD 4 \eTD \bTD 1 \eTD \bTD 1 \eTD \eTR
\bTR \bTD {\bf Total} \eTD \bTD \eTD \bTD \eTD \bTD {\bf 71} \eTD \bTD {\bf 41} \eTD \bTD {\bf 40} \eTD \eTR
\eTABLE
\blank[medium]

El criterio de igualdad es el de la sección \in[sec:objetivo]. En el estudio 11, el único par
resuelto por los dos difiere en el segundo decimal (62.963,65 contra 62.963,64), una diferencia
relativa de $1{,}6 \cdot 10^{-7}$, dentro de la tolerancia.
\stopsubsection

\startsubsection[title={Resultados}]

\startitemize
\item {\bf 40 de 41 pares con el mismo objetivo.} El restante es el punto de 3 lotes en 6×6 del
estudio 08: el monolítico reportó 26.709,36 y el Split 26.709,60, los dos declarando
optimalidad. La diferencia relativa es de $-9{,}0 \cdot 10^{-6}$: supera la tolerancia de
$10^{-6}$, así que con el criterio del comparador sería un \quote{objetivo distinto}.
\item {\bf Mismo estado en los infactibles:} 4 puntos del estudio 08 (2, 3, 4 y 6 lotes en 4×4) son
infactibles en los dos. En el Split, ya en la Etapa 1.
\item {\bf Donde el monolítico no terminó,} en 3 puntos encontró el óptimo del Split sin poder
probarlo: a los 85~s con una brecha de 2,96~\% (estudio 08), a los 109~s con 1,7~\% (estudio 10) y
a los 972~s con 10,37~\% (estudio 11). Son consistentes con la equivalencia, aunque no la prueban.
\stopitemize
\stopsubsection

\startsubsection[title={Las dos diferencias numéricas}]

Hay dos casos en los que el objetivo del monolítico quedó apenas por debajo del óptimo probado del
Split: los tests 13 y 14 ($3{,}83 \cdot 10^{-6}$ absoluto, dentro de la tolerancia) y el punto de 3
lotes en 6×6 ($0{,}24$ absoluto, fuera de ella). Por la demostración de la sección
\in[sec:demostracion], el monolítico no puede tener un óptimo menor que el del Split en
aritmética exacta. La hipótesis registrada es numérica, y no está verificada: con $M = 10^{6}$, un
binario apenas fraccionario dentro de la tolerancia de integralidad de SCIP afloja una
restricción Big-M en cerca de una unidad, y la solución del monolítico queda levemente infactible
sin que SCIP lo detecte.

Una forma de verificarla sería pasar esa solución del monolítico por el verificador del nivel 2 o
revisar la violación máxima de sus restricciones con los valores sin redondear.
\stopsubsection

\stopsection

%%%%%%%%%%%%%%%%%%%%%%%%%%%%%%%%%%%%%%%%%%%%%%%%%%%%%%%%%%%%%%%%%%%%
\startsection[title={Lo que la validación no cubre}, reference=sec:limitaciones]

\startitemize
\item {\bf Ninguna comparación de objetivos en stress.} El monolítico no termina ninguno de los 19
tests de stress, así que en stress solo hay verificación de factibilidad (línea 4), no de
optimalidad. Además, sus stress 05 a 19 nunca se corrieron con $W^{c} = 27$; solo el 04, que
también dio timeout.
\item {\bf Tolerancia de la cadena de frío sin test.} La diferencia entre el filtro exacto del Split
y la tolerancia de SCIP en el monolítico (sección \in[sec:condiciones]) está documentada, pero el
test de tolerancia (35b) solo prueba el lado de \quote{entra}: un par apenas por debajo del límite
sigue elegible en los dos.
\item {\bf Las dos diferencias numéricas,} con una hipótesis sin verificar (sección \in[sec:pares]).
\item {\bf El test 02 no se comparó por separado.} El comparador empareja las filas por nombre de
método, y los tests heredados 02 y 03 tienen el mismo nombre (\ruta{test_two_orders_two_aisles}):
la fila del 02 queda pisada por la del 03. Los dos objetivos del 02 coinciden en los CSV (795 en
los dos modelos), pero sus soluciones no se compararon. Corregirlo es emparejar por archivo y
método.
\item {\bf Sin referencia en el 30a,} por el timeout en los dos modelos, y sin equivalente del 24b en
el monolítico.
\item {\bf La Etapa 1 puede reportar \quote{factible}.} El pipeline del Split acepta ese estado. En los
tests no ocurre, porque la Etapa 1 corre sin límite de tiempo, pero si ocurriera, la condición 3
de la sección \in[sec:condiciones] dejaría de valer.
\item {\bf Resultados anteriores al cambio de nombres.} El 29-09 se unificaron los nombres de
variables y parámetros en el código y en los resultados. Los CSV y los JSON de la comparación se
migraron con el mismo mapeo, sin volver a correr; después del cambio solo se corrieron tests
puntuales de cada modelo, todos sin cambios. Una corrida completa de las dos suites, seguida del
comparador, cerraría ese punto.
\item {\bf El veredicto de los 9 stress del Split no está guardado} en un archivo de resultados
(sección \in[sec:verificador]).
\stopitemize
\stopsection

%%%%%%%%%%%%%%%%%%%%%%%%%%%%%%%%%%%%%%%%%%%%%%%%%%%%%%%%%%%%%%%%%%%%
\startsection[title={Síntesis para la tesina}, reference=sec:sintesis]

\startitemize
\item {\bf Qué se puede afirmar:} el monolítico y el Split son equivalentes en factibilidad y en
valor óptimo. Lo respaldan una demostración en los dos sentidos, válida bajo cuatro condiciones
explícitas, y cuatro líneas empíricas: los 49 pares funcionales resueltos por los dos tienen el
mismo objetivo y los 12 infactibles lo son en los dos; 40 de los 41 pares de la exploración tienen
el mismo objetivo (el restante, a $9 \cdot 10^{-6}$, con una explicación numérica pendiente); y las
9 soluciones de stress del Split son factibles para el monolítico y le cuestan lo mismo.
\item {\bf Qué no se puede afirmar:} que devuelvan las mismas soluciones. En 21 de los 49 pares
funcionales la solución óptima es distinta: rutas del mismo largo, rutas más largas que llegan a
tiempo gracias a la espera voluntaria, o empates de asignación. El Split siempre devuelve la ruta
mínima; el monolítico no.
\item {\bf Cómo presentar la metodología:} conviene contarla como cinco líneas que se complementan
(sección \in[sec:metodologia]), con el criterio de aceptación de cada una, y con dos decisiones
explícitas: la equivalencia se juzga por estado y objetivo, no por solución; y en stress, donde
el monolítico no termina, la pregunta pasa de \quote{¿mismo óptimo?} a \quote{¿la solución del
Split es válida para el monolítico?}.
\item {\bf Por qué importa:} la equivalencia es lo que permite usar al Split, mucho más rápido a
escala, como referencia exacta del problema para las heurísticas, y al monolítico, que es la
formulación directa, como verificador de cualquier solución.
\stopitemize
\stopsection

%%%%%%%%%%%%%%%%%%%%%%%%%%%%%%%%%%%%%%%%%%%%%%%%%%%%%%%%%%%%%%%%%%%%
\startsection[title={Apéndice: archivos y comandos}, reference=sec:apendice]

\startsubsection[title={Dónde está cada cosa en el repositorio de implementación}]
\startitemize
\item {\bf Modelos:} \ruta{models/opp_v5_1_mono/opp_model_v5_1_mono.zpl},
\ruta{models/opp_v5_2_split/opp_model_v5_2_split_stage1_routing.zpl} y
\ruta{models/opp_v5_2_split/opp_model_v5_2_split_stage2_scheduling.zpl}.
\item {\bf Tests:} \ruta{tests/opp_v5_1_mono/} y \ruta{tests/opp_v5_2_split/}; escenarios en
\ruta{tests/shared_scenarios/}; adaptadores de los heredados en
\ruta{tests/infra/factories/migrated_case.py} y \ruta{migrated_mono_case.py}.
\item {\bf Verificación:} \ruta{verification/compare_v5.py} (comparación cruzada),
\ruta{verification/mono_compliance.py} (verificador), \ruta{verification/solution_store.py}
(formato de las soluciones guardadas).
\item {\bf Resultados:} \ruta{results/opp_v5_1_mono/functional/} y
\ruta{results/opp_v5_2_split/functional/} (CSV de nombre fijo y carpetas \ruta{solutions/});
\ruta{results/comparability/compare_v5_split_vs_mono.csv} (veredictos).
\item {\bf Pares de la exploración:} \ruta{experiments/opp_v5/studies/08_mono_frontier/} a
\ruta{11_mono_long_cap/}.
\item {\bf Documento de comparación:}
\ruta{docs/thesis/opp_v5_1_mono_vs_opp_v5_2_split_comparison.tex}.
\stopitemize
\stopsubsection

\startsubsection[title={Comandos}]

Necesitan \ruta{zimpl} en el \ruta{PATH} y PySCIPOpt.

\starttyping
# Suites funcionales con las soluciones guardadas
export WRITE_SOLUTIONS=1
python -m tests.run_suite opp_v5_2_split functional
python -m tests.run_suite opp_v5_1_mono functional

# Comparación cruzada, con la verificación de los stress del Split
python -m verification.compare_v5 --verify-split-stress
\stoptyping
\stopsubsection

\stopsection

\stoptext
